\documentclass[10pt,a4paper]{article}
\usepackage[margin=22mm]{geometry}
\usepackage{amsmath,amssymb,booktabs,array}
\usepackage[colorlinks=true,linkcolor=blue,urlcolor=blue]{hyperref}
\usepackage{microtype}
\setlength{\parindent}{0pt}
\setlength{\parskip}{5pt}
\title{RM9-R1: What does the aggregate observer measure?\\Source findings and a concrete contract decision}
\author{Method-first ramp-metering review}
\date{5 October 2026}
\begin{document}
\maketitle
\textbf{Status.} The accepted source-first review has reached UG4-O.
No speed map, new physical law or simulation is adopted. The new section-speed
Candidate C and the historical passage Candidate A are distinct hypotheses.
Physical passage realization Candidate P keeps the mapping unresolved until
further analytic grounding. Detailed Pros/Cons and the decision are in the
companion findings. This handout supplies the mathematical definitions.

\section{Three different observation supports}
Let $i$ identify a longitudinal road cell, $\ell\in\{1,2\}$ a mainline lane,
$L_i$ its length, $h$ a common time step and $t_s=t_0+sh$.
Stored continuous mass $n_{i\ell}^s$ is measured in vehicle equivalents.
Define road-total mass and road-total density (summed over lanes):
\begin{equation}
N_i^s=n_{i1}^s+n_{i2}^s,\qquad K_i^s=N_i^s/L_i.
\end{equation}
An average density per lane would be $K_i^s/2$; the benchmark uses the summed
normalization consistently. True spatial inventory is an idealized observation,
not detector occupancy or density inferred from detector speed.

At a labeled face $e$, let $q_{e\ell}^s$ be an actual accepted longitudinal
lane rate during the completed interval $[t_s,t_{s+1})$; $Q_e^s$ is its road
total and $B_e^s$ is cumulative passage mass:
\begin{equation}
Q_e^s=\sum_\ell q_{e\ell}^s,\qquad
B_e^{s+1}-B_e^s=hQ_e^s.
\end{equation}
At a merge, the face label must distinguish donor-mainline passage from
combined downstream passage including ramp entry. Counts and interval rates
need not be independent information, but both are accounted for.

\textbf{Passage time mean.} For vehicles $a$ crossing a station/window with
longitudinal speeds $u_a$, the arithmetic passage mean is
\begin{equation}
V_{\rm pass}=\frac{1}{m}\sum_{a=1}^{m}u_a.
\end{equation}
If a lane has one representative passage speed, pooling lanes uses passage
counts (equivalently constant interval rates) as weights. This is not the
speed of the traffic remaining inside an upstream cell. At zero passage the
statistic is unavailable, not a measured zero.

\textbf{Instantaneous section mean.} A density-weighted longitudinal mean
over the traffic present in a section is
\begin{equation}
V_{\rm sec}(t)=
\frac{\sum_\ell\int_{\text{cell }i} k_\ell(x,t)u_\ell(x,t)\,dx}
{\sum_\ell\int_{\text{cell }i} k_\ell(x,t)\,dx},\qquad N_i(t)>0.
\end{equation}
This remains meaningful when vehicles are present but no one crosses the
face. Neither this definition nor the passage definition requires an
upstream-cell-to-face equality.
\newpage

\section{Same-region identities are not arbitrary cell/face identities}
For a rectangular longitudinal space--time region $\mathcal A$ of area
$L\Delta t$, let $d_a$ be the longitudinal distance traveled by vehicle $a$
inside it, and $r_a$ its time spent inside it. The generalized definitions are
\begin{equation}
q_{\mathcal A}=\frac{\sum_a d_a}{L\Delta t},\qquad
k_{\mathcal A}=\frac{\sum_a r_a}{L\Delta t},\qquad
u_{\mathcal A}=\frac{\sum_a d_a}{\sum_a r_a}
=\frac{q_{\mathcal A}}{k_{\mathcal A}}.
\end{equation}
For multiple lanes, sum longitudinal distance and time across the lanes using
the same normalization. Lateral motion itself is not extra longitudinal distance.
These source definitions share the \emph{same region}. Their flow is a traveled
distance average, not necessarily a single cell's downstream accepted rate.
Their density is a time--space average, not necessarily $N_i^s/L_i$ at one
instant. Therefore the identity does not justify defining independent passage
speed as $Q_e^s/K_i^s$ during a transient or receiving-limited interval.

For a section with no internal sources, conservation gives
\begin{equation}
N(t_b)-N(t_a)=B_{\rm in}(t_a,t_b)-B_{\rm out}(t_a,t_b).
\end{equation}
It determines inventory change, not every interior traveled distance or
crossing-speed distribution. Ramp entry is an additional labeled input where
present. Unpassed traffic inside the section must not disappear from estimation.

\textbf{Source outcome.} FB0030 uses one-minute detector flow and speed to
estimate density and acknowledges time/space-mean mismatch. FHWA distinguishes
point passage and section means; dual loops can separately measure crossing
speed. The inspected Jamshidnejad/De Schutter papers define generalized metrics
and retain internal traffic across windows. The 2018 multilane procedure
explicitly excludes lane changes in Section 4.2; a later trajectory approximation
does not validate our adapted lane-changing CTM observer.

\textbf{Evidence boundary.} These are metric and estimation precedents.
No source proves our exact full-history information collision, useful-action
distinction or combined gradient chain. First-order observation closures below
are reviewer proposals, not newly validated field measurements.

\section{Candidate C: section speed under a declared first-order closure}
Let $v_{i\ell}$ be free-flow speed, $w_{i\ell}$ backward-wave speed,
$C_{i\ell}$ capacity and $k_{i\ell}^{\rm jam}=M_{i\ell}/L_i$ jam density.
Use the capacity-limited FD implied by these public parameters:
\begin{equation}
f_{i\ell}^{\rm eq}(k)=\min\{v_{i\ell}k,\ C_{i\ell},\
w_{i\ell}(k_{i\ell}^{\rm jam}-k)\},\qquad
U_{i\ell}^{\rm eq}(k)=f_{i\ell}^{\rm eq}(k)/k\quad(k>0).
\end{equation}
Candidate C assigns this one representative speed to each occupied lane cell
at the stored-state epoch, with $k_{i\ell}^s=n_{i\ell}^s/L_i$:
\begin{equation}
V_i^s=\frac{\sum_\ell n_{i\ell}^s U_{i\ell}^{\rm eq}(k_{i\ell}^s)}{N_i^s},
\quad N_i^s>0.
\label{eq:section}
\end{equation}
An empty lane contributes zero to the numerator without evaluating $0/0$;
an empty cell has an unavailable-speed flag. Occupied stationary traffic can
give $V_i^s=0$. This closure does not add a dynamic speed state or modify G2.
It is a spatially homogeneous/equilibrium observation hypothesis, not a
guarantee about shocks, acceleration or actual detector speeds.
\newpage

\section{What C does and does not hide}
Equation~\eqref{eq:section} equals
$\sum_\ell f_{i\ell}^{\rm eq}(k_{i\ell}^s)/K_i^s$ under the homogeneous
closure. The numerator is a constitutive \emph{within-cell equilibrium} flow
sum, not the actual accepted downstream flow $Q_e^s$. Thus the proposed speed
channel is computed from the lane state separately from accepted face counts;
no statistical or algebraic independence of all observations is claimed.

On the two-lane free-flow FD subcase with distinct known speeds,
\begin{equation}
N_i^s V_i^s=v_{i1}n_{i1}^s+v_{i2}n_{i2}^s,\qquad
n_{i1}^s=N_i^s\frac{V_i^s-v_{i2}}{v_{i1}-v_{i2}},\quad
n_{i2}^s=N_i^s-n_{i1}^s.
\end{equation}
This identifies the \emph{stored} split in that subcase, even if a downstream
constraint caps the face flow, provided the proposed representative-speed
closure still places both cell densities on their free-flow FD branches.
It illustrates the information consequence of a new observation hypothesis;
it is not a measured receiving-limited speed theorem, a full R2 audit or a
matched-history example. Equal public lane speeds and other FD branches need
their own analysis; do not select a symmetry to force a favorable conclusion.

\section{Candidate A and the physical-passage alternative P}
The historical idealized passage proposal uses post-lateral masses
$n_{i\ell}^{L,s}$ and actual lane face flows during the completed step:
\begin{equation}
V_{e,\rm A}^s=
\frac{\sum_\ell q_{e\ell}^s
U_{i\ell}^{\rm eq}(n_{i\ell}^{L,s}/L_i)}{Q_e^s},\qquad Q_e^s>0.
\end{equation}
Besides equilibrium cell speed, this explicitly assumes that the upstream
post-lateral representative speed is the representative crossing speed at
that face. The source review does not validate that extra approximation in
congestion. Zero passage gives an unavailable statistic. A stays viable as
a conditional mathematical benchmark; its operational name is not proof
of physical passage realization.

Candidate P retains actual point passage speed as indispensable, but stops
information testing until a physical/spatial reconstruction law is admitted.
The current inventory/accepted-flow update alone does not provide that law.
Initial conditions, internal queues and any added trajectory approximation
must be reviewed; no microscopic state or simulator is automatically adopted.

\textbf{Recommendation for discussion: C.} It matches the declared cell
inventory's section support and removes the extra cell-to-face assumption.
It remains an explicitly idealized first-order benchmark and may expose lane
contrast. If physical point-detector fidelity is required now, choose P.
Neither the selected source-first route nor this recommendation adopts C/A/P.
\newpage

\section{Causal history and initialization}
Keep the approved ordering
\begin{equation}
n^s\ \xrightarrow{\rm lateral}\ n^{L,s}\
\xrightarrow{\rm forward/merge}\ n^{s+1}.
\end{equation}
The decision epoch $t_*$ occurs at a stored-state epoch after all preceding
steps have completed and before the next update. A future common open-loop
schedule is selected then; this does not introduce repeated feedback control.
With C, observe $(N_i^s,V_i^s)$ at stored epochs $t_s\leq t_*$ and counts/flows
only for completed intervals ending no later than $t_*$. With A, passage speeds
also belong only to those completed intervals. Exclude all future observations
generated by any candidate schedule from the initial information set.

Denote the admitted history by
\begin{equation}
\mathcal H_*=(\mathcal O_{[t_{\rm init},t_*]},\ p,\ b_{\rm public},\ u_{\rm past}).
\end{equation}
Here $\mathcal O$ contains the declared mainline/ramp observations, $p$ the
public model parameters, $b_{\rm public}$ the common known forcing (including
declared future forcing), and $u_{\rm past}$ the prior meter permissions.
Every claimed latent initial state must match every initial and historical
channel. The public graph specifies section lengths, face locations, boundary
lane-allocation rules and downstream supplies. No lane-dependent exogenous
input is kept private merely to create ambiguity.

\textbf{Trivial sufficiency boundary.} If the complete initial physical state
$x^0$, all forcing $b^s$, parameters and previous controls are known, deterministic
recursion $x^{s+1}=F_B(x^s,b^s,u^s)$ already determines the full state. A known
empty initialization is such a subcase when all modeled states are empty and
all structural inputs are public. Hiddenness cannot be claimed there by
ignoring the recursion. If only initial aggregate observations are admitted,
unobserved lane coordinates remain candidates only within their full
observation-consistent set. This is an initialization contract, not permission
to invent private prehistory.

\section{Ramp balances and geometry requirements}
Let $W_r^s$ be external ramp waiting, $R_{\rm pre}^s$ traffic upstream of the
meter inside finite ramp storage, and $R_{\rm post}^s$ traffic in any declared
post-meter region. Let $A_r^s$ be external arrival rate, $b_r^s$ admitted entry,
$g_r^s$ actual meter release and $r_r^s$ actual merge entry, all in mass/time.
Where these compartments exist, conservation requires
\begin{align}
W_r^{s+1}&=W_r^s+h(A_r^s-b_r^s),\\
R_{\rm pre}^{s+1}&=R_{\rm pre}^s+h(b_r^s-g_r^s),\\
R_{\rm post}^{s+1}&=R_{\rm post}^s+h(g_r^s-r_r^s).
\end{align}
For several post-meter cells, the last line is their sum balance and internal
flows cancel. The permission limits $g_r^s$; finite sending/storage/receiving
can further restrict it. $r_r^s$ is set by the adopted shared receiver using
the downstream ramp offer, not automatically by the current permission.
\newpage

\section{Minimum layout/accounting declaration before R2/G3}
Declare a finite two-mainline-lane cell graph, the meter location, any finite
pre/post-meter ramp regions, the ramp-entry receiver and every observed
section/face. Before and after the merge, donor passage and combined passage
must have distinct labels. At a simple shared receiver,
\begin{equation}
q_{\rm combined}^s=m^s+r_r^s
\end{equation}
for its two inputs; a second through lane is accounted for separately in the
road-total downstream face. Ordinary through-lane transfers and ramp entry
must jointly respect each receiving lane's approved supply. This is a required
interface declaration, not adoption of a new topology or mandatory-merge law.

If the meter and merge are truly co-located, omit nonexistent post-meter storage
and specify that specialization explicitly. Do not impose co-location merely
to identify release and entry. Conversely, storage/propagation cannot be
silently omitted where the task includes them. Numerical cell counts are not
a separate methodological gate, but the geometry must be coherent before
a realized full-history/action certificate.

The complete population includes mainline external waiting where present,
mainline cell mass, $W_r$, $R_{\rm pre}$ and $R_{\rm post}$. All initial traffic
and unserved demand remain counted. The accepted total-time objective and
UG1--UG3 dynamics are unchanged. No gradient experiment was run.

\section{Evidence and disposition}
\begin{tabular}{>{\raggedright\arraybackslash}p{0.26\textwidth}>{\raggedright\arraybackslash}p{0.66\textwidth}}
\toprule
Established in R1 & Source-backed metric/support distinctions; conditional
observer proposals; causal/ramp/initialization requirements.\\
Still open & UG4-O observer adoption; source-grounded concrete layout; G3
feasible equal-history/useful-action evidence; G4 validated useful gradients.\\
Not authorized & Traffic runs, SUMO installation, new dynamics, extra lanes,
noise/smoothing, successor or deferred extension execution.\\
\bottomrule
\end{tabular}

\textbf{Primary references inspected.}
\begin{itemize}
\item \href{https://repository.tudelft.nl/file/File_2c4359e7-5e5a-4651-964f-6f3610ab89ae}{FB0030: lane model}, printed 1768--1769.
\item \href{https://www.fhwa.dot.gov/policyinformation/tmguide/tmg_2022/traffic-data-collection.cfm}{FHWA Traffic Monitoring Guide 2022}, speed and occupancy sections.
\item \href{https://ops.fhwa.dot.gov/publications/fhwahop08038/04infoshare.htm}{FHWA Data Quality White Paper}, Section 4.3.2.
\item \href{https://www.dcsc.tudelft.nl/~bdeschutter/pub/rep/15_035.pdf}{Jamshidnejad/De Schutter 2015}, Sections I--II.
\item \href{https://www.dcsc.tudelft.nl/~bdeschutter/pub/rep/17_001.pdf}{Jamshidnejad/De Schutter 2018 author report}, Sections 2.1, 4.1--4.2 and 5.
\end{itemize}
\textbf{Stop.} Discuss C/A/P at UG4-O. No observer is adopted by this document.
\end{document}
